\section{Related work}
\label{sec:related}
\paragraph{Multi-task benchmarks as voting.}
\citet{zhang2024tradeoffs} and \citet{lanctot2023vase} view tasks as voters and models as candidates; the former separates \emph{cardinal} benchmarks (aggregating scores) from \emph{ordinal} ones (aggregating per-task rankings), applies Arrow's theorem to ordinal ones and measures sensitivity to irrelevant changes; the latter proposes voting rules as evaluation frameworks. These works ask how a rule responds to changes in the set of models or tasks; we ask what it costs a \emph{single} model to win. \citet{dehghani2021lottery} show that the choice of tasks can reorder methods.
\paragraph{Manipulation in social choice.}
Minimum-cost election manipulation is classical in computational social choice (bribery, shift bribery, margin of victory; \citealp{faliszewski2006bribery}; \citealp{brandt2016handbook}). That literature is ordinal and studies computational complexity. Ours is the cardinal counterpart: the target's task scores move continuously inside a box, rivals are fixed, and the rules are score aggregators (mean, median, minimum, generalized means) and, for comparison, ordinal rules (Borda, plurality, a Condorcet requirement). The closed forms for the mean, minimum and odd median are elementary (they are the cardinal analogue of a margin-of-victory computation), so what we add is the exact computation on real score tables for every model, the reversals of the rule ordering, and the dependence on cap and frame.
\paragraph{Strategic behaviour and contamination.}
\citet{chen2026incentives} model benchmarking as a designer--developer game: each developer has a latent capability, a scalar score and a scalar effort, and is paid by rank. Any pure equilibrium keeps the capability order, the only candidate is zero effort, and it exists iff the cost of the \emph{just-overtake} effort (to pass the model directly above) is at least the adjacent reward gap (their Prop.~4.5, Thm.~4.6); under winner-take-all the binding pair is second against first. Their overtaking cost is the single-score counterpart of our cost-to-first, and tune-before-test raises it. Multiple tasks enter only through one fixed linear weight on the latent capability, with no per-task effort and no choice of aggregation rule, so their results order no rules; the rule, task-by-task and frame dependence measured here lie outside them. Their evidence is a case study of five Qwen2.5 sizes; ours is every model of two leaderboards under fixed rivals, with no equilibrium claim. \citet{singh2025illusion} document distortions of the Chatbot Arena leaderboard (private testing, sampling). \citet{zhou2023cheater} and \citet{yang2023rephrased} show that training on benchmark-related data inflates scores and can evade decontamination. These motivate our investment model but concern measured effects of leakage; we do not measure achievable gains and treat the per-task gain cap as a scenario.
\paragraph{Benchmarks used.}
We use per-subject MMLU \citep{hendrycks2020mmlu} accuracies from HELM \citep{liang2022helm} and the six tasks of the Open LLM Leaderboard (v2; IFEval, BBH \citealp{suzgun2022bbh}, MATH Lvl 5, GPQA, MuSR, MMLU-Pro), whose scores we take as published.
